\documentclass[11pt]{article}
\usepackage[margin=1in]{geometry}
\usepackage{amsmath,amssymb}
\usepackage{booktabs}
\usepackage{array}
\usepackage{hyperref}
\hypersetup{colorlinks=true,linkcolor=black,citecolor=black,urlcolor=blue}

\title{\textbf{Model Training Results}\\[2pt]\large Detection and Attribution of AI-Generated and Human Written
Phishing Emails from Legitimate Emails}
\author{Mannan Bhardwaj -- 25BYB0013 \\ Bhoumik Singh Rajput -- 25BYB0120 \\[4pt]
School of Computer Science Engineering and Information Systems}
\date{}

\begin{document}
\maketitle

\section{Overview}
This is the first run of the full pipeline end to end: real feature extraction (not proxy features),
real statistical gating on the extracted values, parameter locking, and evaluation on the independent
holdout set. Two findings from this run contradict earlier, proxy-feature-based claims and are reported
here rather than quietly corrected.

\section{Feature Extraction -- what was actually used}
All 20 features were computed with real tooling, not approximations: spaCy (\texttt{en\_core\_web\_sm})
for dependency parsing and POS tagging, a from-scratch trigram language model (stupid backoff) trained
once on the training set's own text for token-predictability features, VADER for sentence-level
sentiment, and \texttt{pyspellchecker} for orthographic error detection. A pretrained transformer
language model was considered for the token log-probability feature and rejected: it was unreachable
from this environment (no route to download pretrained weights), and a heavy black-box model would also
cut against the project's own stated design goal of a lightweight, fully auditable pipeline. Email text
is capped at 6{,}000 characters before parser/tokenizer-based features run, since a handful of extreme
outliers (up to 31{,}396 words, almost certainly mbox-thread artifacts) would otherwise dominate runtime
for no signal benefit.

\section{Statistical Gating -- real results}
Kruskal-Wallis, Mann-Whitney U, Chi-Square, and joint Benjamini-Hochberg FDR correction were re-run on
the real extracted values (Section~3 methodology, unchanged). Table~\ref{tab:gate} summarizes the
outcome.

\renewcommand{\arraystretch}{1.25}
\begin{table}[h]
\centering
\small
\begin{tabular}{@{}lrrcrrc@{}}
\toprule
& \multicolumn{3}{c}{\textbf{Main (3-group)}} & \multicolumn{3}{c}{\textbf{Attribution (2-group)}} \\
\textbf{Feature} & adj.\ $p$ & effect & gate & adj.\ $p$ & effect & gate \\
\midrule
avg\_token\_logprob            & $\approx\!0$ & 0.956 & \checkmark & $\approx\!0$ & $-$0.956 & \checkmark \\
sentence\_perplexity\_variance & $\approx\!0$ & 1.251 & \checkmark & $\approx\!0$ & 0.835    & \checkmark \\
sentence\_length\_std          & $\approx\!0$ & 0.630 & \checkmark & $\approx\!0$ & 0.263    & \checkmark \\
max\_parse\_depth              & $\approx\!0$ & 0.637 & \checkmark & $\approx\!0$ & $-$0.637 & \checkmark \\
paragraph\_density             & $\approx\!0$ & 3.150 & \checkmark & $\approx\!0$ & $-$3.150 & \checkmark \\
clause\_density                & $\approx\!0$ & 0.405 & \checkmark & $\approx\!0$ & $-$0.405 & \checkmark \\
orthographic\_error\_density   & $\approx\!0$ & 1.431 & \checkmark & 0.0099 & 0.406 & \checkmark \\
punctuation\_entropy           & $\approx\!0$ & 0.889 & \checkmark & $\approx\!0$ & $-$0.541 & \checkmark \\
urgency\_density               & $\approx\!0$ & 1.325 & \checkmark & $\approx\!0$ & $-$0.042 & fails \\
pronoun\_balance               & $\approx\!0$ & 0.992 & \checkmark & $\approx\!0$ & $-$0.044 & fails \\
sentiment\_variance            & $\approx\!0$ & 0.713 & \checkmark & $\approx\!0$ & $-$0.631 & \checkmark \\
imperative\_density            & $\approx\!0$ & 0.225 & \checkmark & $\approx\!0$ & 0.225    & \checkmark \\
lexical\_diversity             & $\approx\!0$ & 0.514 & \checkmark & $\approx\!0$ & $-$0.395 & \checkmark \\
pronoun\_density               & $\approx\!0$ & 0.969 & \checkmark & $\approx\!0$ & $-$0.354 & \checkmark \\
generic\_greeting\_flag        & $\approx\!0$ & 0.411 & \checkmark & $\approx\!0$ & 0.240 & \checkmark \\
subject\_refw\_flag            & $\approx\!0$ & 0.511 & \checkmark & $\approx\!0$ & 0.158 & \checkmark \\
attachment\_mention\_flag      & $\approx\!0$ & 0.168 & \checkmark & $\approx\!0$ & 0.120 & \checkmark \\
bracketed\_cta\_flag           & $\approx\!0$ & 0.377 & \checkmark & 0.0004 & 0.065 & fails \\
sensitive\_info\_flag          & $\approx\!0$ & 0.306 & \checkmark & $\approx\!0$ & 0.175 & \checkmark \\
threat\_clause\_flag           & $\approx\!0$ & 0.168 & \checkmark & 0.0055 & 0.051 & fails \\
\bottomrule
\end{tabular}
\caption{Dual-gate results on real extracted features (adjusted $p<0.05$ \textbf{and} effect size at or
above threshold: $|d|\geq0.20$ continuous, $V\geq0.10$ categorical). Effect-size sign for the attribution
column indicates direction (which class scores higher), not magnitude -- the gate itself is applied to
the absolute value.}
\label{tab:gate}
\end{table}

\textbf{All 20 of 20 features clear the main gate}; \textbf{16 of 20 clear the attribution gate}
(urgency\_density, pronoun\_balance, bracketed\_cta\_flag, and threat\_clause\_flag do not -- each is
statistically significant purely from sample size, but too small an effect to usefully separate
Human-Phishing from AI-Phishing).

\subsection*{A claim this contradicts}
\texttt{proposal.tex} currently describes imperative verb density as ``the single strongest attribution
signal in our feature set.'' Measured against real extracted data, it is not: its attribution effect size
($d=0.225$) is one of the smallest that still clears the gate, while \texttt{paragraph\_density}
($d=3.15$), \texttt{avg\_token\_logprob} ($d=0.96$), and \texttt{sentence\_perplexity\_variance}
($d=0.84$) are all substantially larger. The proposal's claim was based on an earlier qualitative
expectation, not this measurement, and should be corrected before the paper is finalized.

\subsection*{A result worth treating with suspicion, not celebration}
\texttt{paragraph\_density}'s effect size ($d=3.15$) is far larger than any other feature's -- large
enough to be a red flag rather than good news. The three classes come from three different raw
sources with three different formatting conventions (Nazario phishing text, raw Enron plain-text email,
and LLM-generated output), so a feature this strongly separated is more likely picking up a
corpus-formatting artifact (how each source happens to insert paragraph breaks) than a genuine
linguistic signal that would generalize to phishing emails collected a different way. This is the same
category of caveat already on record for other features (the corpus-era markers in the Legitimate class,
the AI-Phishing corpus's 0\% exclamation marks). It is flagged here rather than presented as a strong
result.

\section{Parameter Locking -- two pipelines, not one}
A design decision was required that the proposal's equations did not fully pin down: whether Step 2
(attribution) reuses the same 3-class posterior as Step 1, or is a separately trained model restricted to
phishing-only data and the attribution-gate feature subset. The abstract's own wording -- ``a dedicated
attribution-specific feature subset'' -- only makes sense as a separate model, since several main-gate
features (urgency\_density, pronoun\_balance) do not actually hold up for attribution. \textbf{Two
pipelines were therefore built}: a 3-class \emph{main pipeline} (Legitimate / Human-Phishing /
AI-Phishing, 20 features) for Step 1, and a 2-class \emph{attribution pipeline} (Human-Phishing /
AI-Phishing, trained only on phishing rows, 16 features) for Step 2.

Two other open parameters were fixed: \textbf{8 quantile bins} for the Rank-Binned NB (a middle ground
for roughly 1{,}497 training rows per class), and the Meta-NB combination was implemented as
$P_{\text{meta}}(C_k) \propto P_1(C_k)\cdot P_2(C_k) \,/\, P(C_k)$, renormalized over $k$ -- the
self-consistent form of Eq.~9 once $P(P_m\mid C_k)$ is read as $P_m(C_k)/P(C_k)$, which avoids
triple-counting the class prior (each sub-model's own posterior already carries one copy of it).

\section{Evaluation on the Holdout Set}
This is the real generalization test: the holdout set's Legitimate rows are overwhelmingly raw Enron
email (4{,}934 of 4{,}976), not the Nazario-legit mix the model trained on -- deliberately the harder
test. Human-Phishing holdout rows are few (42), a direct consequence of the earlier balancing step, so
that side of the attribution table is noisier than the AI-Phishing side.

\renewcommand{\arraystretch}{1.4}
\begin{table}[h]
\centering
\small
\begin{tabular}{@{}p{4.0cm}p{1.8cm}p{2.9cm}p{5.6cm}@{}}
\toprule
\textbf{Stage} & \textbf{Accuracy} & \textbf{Precision} & \textbf{Recall} \\
\midrule
Step 1 (Legitimate vs.\ Phishing) & 0.854 & 0.776 (Phishing) & 0.915 (Phishing) \\
Step 2 (Attribution, true-phishing rows only) & 0.865 & -- & 0.952 (Human-Phishing, $n{=}42$) / 0.863 (AI-Phishing, $n{=}3489$) \\
End-to-end 3-class & 0.817 & -- & -- \\
\bottomrule
\end{tabular}
\caption{Holdout evaluation. Step 2 is scored by running the attribution model directly on every
true-phishing row, independent of whether Step 1 actually flagged that row as phishing, so Step 1's
errors cannot contaminate this number.}
\end{table}

Step 1's false-positive rate on Legitimate email is 18.9\% (934 of 4{,}934 raw Enron emails flagged as
phishing) -- notably worse than the 99.4\% accuracy of the earlier TF-IDF + Logistic Regression sanity
baseline from the data-validation stage. This gap is expected and informative, not a bug: the baseline
sees the full vocabulary of every email, while this pipeline deliberately restricts itself to 20
interpretable statistical features. The gap is the real, measured cost of that interpretability
constraint, and should be reported as such in the paper rather than compared directly against the
baseline's number without this context.

\section{Open Items}
\begin{enumerate}
\item \texttt{proposal.tex} needs correcting on the imperative-density claim above, and updating from
14 to 20 features (Section 2.1), which was already outstanding.
\item The large \texttt{paragraph\_density} effect size should be investigated further before being
reported as a finding -- specifically, whether it survives when tested only within a single raw source
per class rather than across sources.
\item The two-pipeline architecture decision (Section~4 above) should be reviewed and either confirmed
or revised, since it is this report's reading of an ambiguous part of the proposal, not something the
proposal stated outright.
\end{enumerate}

\end{document}
